\documentclass[../main.tex]{subfiles}
\begin{document}
%==========================================TIÊU ĐỀ TẬP========================================
\begin{table}[h]
  \begin{adjustwidth}{-1cm}{}
    \begin{tabular}{>{\centering\arraybackslash}p{6cm}|>{\raggedright\arraybackslash}p{12.5cm}}
      \multirow{5}{6cm}
      % Cột bên trái: ÉPISODE và số tập
      {\\[-40pt]
        \flaregothic\fontsize{20pt}{20pt}\selectfont \centering
        \color{eptype}ÉPISODE\color{black}\\
        \burbank\fontsize{72pt}{72pt}\selectfont
        \color{epnum}22\color{black}\\[6pt]
        \cafeteria\fontsize{20pt}{20pt}\selectfont\color{epnum}(2-5)\color{black}
        \\[18pt]
        \flaregothic\fontsize{18pt}{18pt}\selectfont
        \color{epattr}ÉPISODE PERDU
        \color{black}
      }
       &
      % Dòng thứ nhất: tên tập tiếng Nhật
      {\fotskip\fontsize{18pt}{18pt}\selectfont \ruby{結}{けっ}\ruby{果}{か}、\ruby{被}{ひ}\ruby{害}{agai}\ruby{総}{そう}\ruby{額}{がく}〇〇\ruby{億}{おく}\ruby{円}{えん}
      } \\[6pt]
      % Dòng thứ hai: tên tập tiếng Anh
       & {\cafeteria\fontsize{18pt}{18pt}\selectfont Damage Reports Billions}                                                                    \\[4pt]
      % Dòng thứ ba: tên tập tiếng Việt
       & {\vnmsans\fontsize{18pt}{18pt}\selectfont Cú xoay người tỷ yên}                                                      \\[8pt]
      % Dòng thứ tư: Nhân vật xuất hiện
       & {\caviar{\fontsize{14pt}{14pt}\selectfont Nhân vật: \emp{kuon1} \emp{shion} \emp{roboco}}} \\[8pt]
      % Dòng cuối cùng: Thời gian diễn ra của tập (thường sẽ là ngày phát hành)
       & {\continuum\fontsize{16pt}{16pt}\selectfont Ngày phát hành: 06/10/2019}                                                                 \\[18pt]
    \end{tabular}
  \end{adjustwidth}
\end{table}
%=========================================TRUYỆN CHÍNH========================================
\color{black}

\txtbx[2]{{\color{vol} \uline{Tóm tắt:}} Chỉ vì một ý tưởng giữ ấm bụng đơn giản cho Shion, Roboco đã vô tình biến nàng Phù thủy thành một mũi khoan sống xuyên thủng tòa nhà văn phòng. Thảm họa tỷ yên này đã khiến Kuon nổi trận lôi đình và buộc cả hai phải tham gia vào công cuộc tái thiết đầy gian nan.}

Một ngày đầu tháng Mười, giữa tiết trời thu se lạnh. Cơn nắng gắt mùa hạ đã dần dịu đi, nhường chỗ cho những cơn gió nhẹ len qua khung cửa, mang theo chút hơi lạnh khiến ai nấy cũng phải cân nhắc khoác thêm chiếc áo mỏng.

Tuy nhiên, trong một góc nhỏ tại văn phòng \hll, nàng Phù thủy trẻ tuổi Shion lại đang đứng ôm bụng với vẻ mặt cực kỳ khổ sở.

``Đau quá\dots\ em đau bụng quá đi mất\dots'' Shion rên rỉ, gương mặt nhỏ nhắn nhăn nhó vì khó chịu.

Tôi lúc này đang chuẩn bị danh sách công việc cho ngày hôm nay, ngẩng đầu lên nhìn. Thấy điệu bộ của em ấy, tôi không khỏi lo lắng, nhưng phần nhiều là cảm giác bất lực. Tôi thở dài một tiếng, đặt cây bút xuống bàn: ``Em đã kêu ca như vậy suốt mười lăm phút rồi đấy. Rốt cuộc là đau kiểu gì mà dai dẳng thế?''

\img{2-6a}

Ở gần đó, Roboco cũng tiến lại gần, cúi xuống nhìn nàng Phù thủy đang run rẩy với vẻ tò mò. Cô nàng hỏi bằng giọng đều đều đặc trưng của người máy: ``Em bị cảm lạnh sao Shion-chan?''

Shion khẽ lắc đầu, đáp lại đàn chị bằng giọng yếu ớt, như thể không còn chút sức lực nào: ``Em cũng không biết nữa\dots\ tự nhiên nó cứ nhói lên thôi\dots''

Tôi khoanh tay, đưa mắt quan sát Shion từ đầu đến chân một lượt, rồi sau đó khẳng định chắc nịch: ``Anh nghĩ chắc chắn là tại bộ trang phục của em đấy. Phơi hết cả bụng ra thế kia giữa cái tiết trời này thì không lạnh bụng mới là lạ. Đau là đúng rồi, chẳng oan đâu.''

Roboco gật gù tán thành ý kiến của tôi: ``Cũng có thể lắm chứ\~{} Thu sang rồi mà mặc hở hang quá là không tốt đâu.''

Nàng phù thủy tóc bạc trắng run lẩy bẩy, hai tay ôm bụng chặt hơn, đôi mắt ngấn lệ đầy vẻ bất lực: ``Nhưng bộ đồ của em vốn dĩ nó đã thiết kế như thế rồi, em biết làm thế nào được bây giờ\dots?''

Roboco khẽ cười, đưa tay lên che môi như đang cố tình trêu chọc: ``Chị đã bảo mà, kiểu gì mặc thế cũng có ngày dính chưởng.''

Tôi cũng tặc lưỡi phụ họa: ``Lần đầu tiên anh thấy Robo-PON nói một câu nghe có lý đấy.''

Nghe thấy biệt danh đó, cô nàng Người máy lập tức quay sang tôi phản đối, vẻ mặt không thoải mái chút nào: ``Em không có PON mà! Đã bảo đừng gọi em như thế rồi!''

Trong khi đó, Shion dù đang đau vẫn cố gân cổ lên cãi lại đàn chị: ``C-Chị đang định khịa em hay gì đấy hả!''

Tôi chỉ biết lắc đầu trước hai cô nàng này: ``Người ta nói đúng quá rồi còn kêu ca cái gì nữa\dots''

Bất chợt, Roboco xoay nhẹ một vòng, bộ trang phục kim loại của cô phản chiếu ánh nắng lấp lánh khiến chúng tôi không khỏi nheo mắt. Với vẻ mặt đầy kiêu ngạo, cô nàng tuyên bố đầy đắc thắng: ``Kuon-senpai nói đúng đó nha\~{}! Dù sao chị cũng là rô-bốt cao cấp, có bao giờ biết lạnh bụng là gì đâu.''

Tôi cười nhạt, phẩy tay như muốn đuổi khéo: ``Rồi rồi, em thì nhất rồi. Rô-bốt mình đồng da sắt, không biết đau, anh thừa hiểu.''

Shion lúc này càng run rẩy dữ dội hơn, cả người co rụt lại như đang cố chống chọi với cái lạnh đang xuyên qua lớp vải mỏng manh. Nhìn nàng Phù thủy khổ sở đến tội nghiệp, tôi quay sang bảo Roboco: ``Nói thế thôi, chứ em xem có cách nào giúp em ấy đi. Mau tìm cái gì đó che cái bụng lại kẻo em ấy ngất ra đây bây giờ.''

Thấy Shion đáng thương quá, Roboco thở dài rồi gật đầu đồng ý. Cô tháo chiếc áo khoác mỏng trên vai mình, nhẹ nhàng buộc quanh vòng eo của Shion như một chiếc đai giữ ấm tạm thời. ``Đây, cho em mượn cái áo khoác của chị nè.''

\img{2-6b}

Shion ngần ngại ngước lên nhìn đàn chị với đôi mắt lấp lánh sự cảm kích: ``Có được không vậy chị? Roboco-senpai không thấy lạnh sao?''

Tôi đứng bên cạnh mỉm cười, giọng nói cũng dịu đi đôi chút: ``Có đồ giữ ấm là tốt rồi. Nhớ tí nữa cảm ơn rồi trả lại cho em ấy đàng hoàng là được.''

``\dots Cảm ơn chị, Roboco-senpai,'' Shion gật đầu, nét mặt dần giãn ra. Cảm giác ấm áp lan tỏa khiến cô nàng rên lên vì sung sướng: ``Ôi chao\~{} ấm quá đi mất\dots\ sống lại rồi\dots''

Tôi khoanh tay, nghiêm nghị dặn dò: ``Đấy, bây giờ bắt đầu vào mùa lạnh rồi, làm ơn nhớ giữ kín cơ thể một chút. Đừng có để thời trang phang thời tiết nữa.''

Shion cười gượng gật đầu, trong khi Roboco đứng bên cạnh trông có vẻ rất tự hào vì hành động nghĩa hiệp của mình.

\ct

Sau một hồi im lặng để Shion tận hưởng hơi ấm, bỗng nhiên đôi mắt của Roboco lóe lên một tia sáng kỳ lạ, như thể bộ não điện tử của cô vừa nảy ra một ý tưởng vô cùng ``đột phá''.

``Này, em có muốn thử làm chuyện này không Shion-chan?''

Shion ngơ ngác nhìn đàn chị, trên đầu hiện ra một dấu hỏi to đùng. Tôi nghe vậy liền nheo mắt cảnh giác: ``Này, anh nhắc trước là đừng có làm mấy cái trò xằng bậy đấy nhé. Văn phòng không phải chỗ chơi đâu.''

\img{2-6c}

Roboco cười tươi rói, chẳng thèm mảy may quan tâm đến lời cảnh báo của tôi: ``Em nghĩ Kuon-senpai lại hiểu nhầm đi đâu rồi. Ý em là muốn cùng Shion-chan thử `xoay vòng\~{}' ấy mà.''

``À, ý em là chơi trò con quay à?'' Tôi thở phào nhẹ nhõm khi thấy ý tưởng có vẻ vô hại. ``Được thôi, nhưng mà tại sao lại muốn chơi trò đó lúc này?''

Roboco chu môi đáp: ``Thì đúng là như thế mà! Chứ anh đang nghĩ cái gì trong đầu vậy hả\~{}?''

Tôi vội xua tay phân tán sự chú ý: ``Không, không có gì. Anh chỉ lo xa chút thôi.''

Shion nghe đến trò xoay vòng thì lập tức hào hứng hẳn lên: ``Ồ, ý hay đấy chị! Nghe có vẻ vui nha!''

``Thế để anh ra ngoài kiếm cái ghế xoay cho mấy đứa--'' Tôi vừa định đứng dậy thì Roboco đã nhanh chóng ngăn lại.

``Không cần đâu anh, em có thể tận dụng luôn cái áo khoác này nè!''

Tôi ngẩn người: ``Ý em là cái áo mà Shion-chan đang buộc quanh bụng ấy hả?''

``Vâng! Chính nó! Shion-chan cũng muốn thử cảm giác mạnh mà, đúng không nào?''

Tôi gãi đầu, cảm thấy có một sự bất an len lỏi trong lòng về cái ý tưởng ``xoay vòng'' thủ công này. Nhưng nhìn ánh mắt đầy phấn khích của hai cô nàng, tôi đành tặc lưỡi chấp nhận: ``Thôi được rồi, dù sao cách này cũng đơn giản. Nhưng Shion-chan chắc chắn là ổn chứ?''

Nàng Phù thủy cười tự tin: ``Em sẽ ổn mà! Em có phép thuật hộ thân cơ mà, mấy cái xoay này có là gì đâu!''

Cuối cùng, tôi đành đứng sang một bên để hai người họ thực hiện, dù trong thâm tâm vẫn không khỏi lo lắng về sự vụng về tiềm ẩn của Roboco. ``\textit{Để xem cô nàng Robo-PON này có làm gì điên rồ không đây\dots}'' Tôi thầm nghĩ, ngồi từ xa quan sát.

\img{2-6d}

Roboco nắm lấy hai vạt tay của chiếc áo khoác đang buộc quanh người Shion, ra hiệu đầy dứt khoát: ``Chuẩn bị nhé, chúng ta bắt đầu đây!''

Shion gật đầu lia lịa: ``Làm đi chị! Em sẵn sàng rồi!''

Với vẻ mặt đầy quyết tâm, cô nàng Người máy đứng vào tư thế chuẩn bị, hai chân bám trụ vững chắc, đôi mắt bỗng chuyển sang màu tím rực rỡ như thể đang kích hoạt một bộ xử lý sức mạnh phi thường nào đó.

\begin{figure}[H]
  \centering
  \begin{subfigure}[b]{0.45\textwidth}
    \centering
    \includegraphics[width=8cm]{../../../../Image/Book2/2-6e1.png}
  \end{subfigure}
  \quad
  \begin{subfigure}[b]{0.45\textwidth}
    \centering
    \includegraphics[width=8cm]{../../../../Image/Book2/2-6e2.png}
  \end{subfigure}
\end{figure}

Tôi nheo mắt nhìn, không khỏi thắc mắc: ``Này\dots\ có nhất thiết phải gồng mình lên kinh khủng như thế không? Anh thấy hơi quan ngại rồi đấy.''

Không buồn trả lời, Roboco dùng hết sức bình sinh giật mạnh tay áo, biến Shion thành một con quay sống động giữa văn phòng.

\gif{2-6f1}{1}{32}

Shion vừa xoay vừa kêu la trong phấn khích: ``Oa\~{}! Cả thế giới đang xoay tròn tròn nè\~{}! Vui quá đi!''

Tôi vội đứng bật dậy: ``Anh nghĩ là em làm hơi quá tay rồi đấy Roboco! Dừng lại mau!''

Nhưng dường như sức mạnh của nàng Rô-bốt là quá lớn, cú xoay mất kiểm soát khiến Shion biến thành một mũi khoan sống với vận tốc kinh hoàng. Em ấy xuyên thủng cả tầng lầu, lao thẳng xuống mặt đất với một tiếng động chói tai.

\gif{2-6f2}{1}{148}

Căn phòng rung chuyển dữ dội như gặp động đất cấp độ cao, mọi thứ xung quanh bắt đầu sụp đổ. Bụi mù mịt bao trùm toàn bộ hiện trường. Một lúc sau, khi đám bụi dần lắng xuống, tôi nhận ra mình vừa thoát chết trong gang tấc. Đứng bần thần trước đống đổ nát tan hoang của văn phòng, tôi bất giác hét lớn bằng cả vốn liếng ngoại ngữ của mình:

``\textbf{{\color{red} C-C-C-Cái củ quặc gì vừa diễn ra thế này!?!!} {\color{blue} Roboco, em vừa làm cái trò mẹ gì vậy hả!?!}}''

Chưa hả cơn giận, tôi gằn giọng tiếp tục: ``\textbf{Em nghĩ mình đang làm cái trò mèo gì thế hả? Định phá nát cả cái công ty này cho sập mới vừa lòng à!?}''

Roboco từ trên không dần hạ cánh xuống mặt đất, gương mặt đỏ bừng, lắp bắp phân bua trong hoảng loạn: ``Ơ\dots\ em không ngờ sức mạnh của mình lại khủng khiếp đến thế\dots\ Em xin lỗi, Kuon-senpai\dots\ em không cố ý đâu\dots''

Dưới mặt đất, Shion bị chôn gần hết người trong đống gạch đá, chỉ còn trơ lại mỗi cái đầu lộ ra ngoài. Có vẻ như cô nàng cũng bị choáng váng sau cú xoay tử thần đó. Thế nhưng, thật kỳ lạ, em ấy lại nở một nụ cười thỏa mãn đến khó hiểu.

\img{2-6g}

``Ôi\dots\ không ngờ là ở sâu dưới lòng đất này lại ấm áp đến thế nhe\dots'' Shion thào thào, gương mặt giãn ra đầy thư thái dù đang bị vùi lấp giữa lửa và khói bụi. Chiếc mũ phù thủy từ đâu bay tới, đáp nhẹ nhàng lên mặt cô nàng, khiến Shion cứ thế nằm im lìm như đang chìm vào một giấc ngủ ngon lành.

Tôi quay phắt lại, gào lên trong sự bất lực tột độ: ``\textbf{Ấm cái khỉ gió ấy!! Tỉnh lại cho anh nhờ cái!!}''

Thật đúng là sự yên bình trước cơn bão. Chỉ từ một ý tưởng giữ ấm đơn giản, hai cô nàng này đã biến nó thành một thảm họa trị giá hàng tỷ yên. 

Tôi chỉ biết ôm trán, nuốt nước mắt vào trong khi nghĩ đến đống hóa đơn sửa chữa: ``\textit{Thế này thì chắc phải xin trợ cấp toàn công ty mới đủ mất\dots\ Giờ thì ăn nói làm sao với A-chan và Tanigo-san đây trời\dots}''

%==========================================TRUYỆN PHỤ=========================================
\omk

Sau khi nàng Phù thủy đã thiếp đi được một lúc, Roboco lén lút tiến lại gần, tay vẫn còn cầm vạt áo khoác gây họa. Thấy mặt tôi đang đỏ gắt lên vì tức giận, cô nàng cười gượng gạo, hỏi nhỏ: ``Ê-tồ\dots\ để em giúp dọn dẹp đống này luôn nhe?''

Thế nhưng, lời đề nghị còn chưa dứt, một chiếc cờ lê còn sót lại từ đống đổ nát đã bay vút tới, găm thẳng vào trán nàng Người máy.

``\textbf{Cái này là để em nhớ mà sống cho biết điều hơn đấy, đồ Robo-PON phá hoại!}'' Tôi hét lên, thực sự đã đạt đến giới hạn chịu đựng.

Cũng chẳng thể trách tôi được, vì đây là lần đầu tiên có người trong \hll\ chơi lớn đến mức làm sập cả tòa nhà văn phòng. Mà kẻ gây chuyện lại là cô nàng vốn đã nổi danh với biệt tài ``đụng đâu hỏng đó'' như Roboco, nên việc tôi nổi điên cũng là chuyện dễ hiểu. Có lẽ đây là lần đầu tiên tôi thực sự giận dữ đến mức mất kiểm soát như vậy.

Roboco ôm trán rên rỉ, nước mắt lưng tròng: ``Đau quá! Anh làm cái gì thế Kuon-senpai!? Sao lại bạo lực với rô-bốt thế này!''

Tôi vẫn không hề hạ hỏa, quát ngược lại: ``\textbf{Đau à!? Thế cái nỗi đau đó có bằng cái đống nợ tỷ yên mà em vừa gây ra không hả!?}''

Nói đoạn, tôi bắt đầu bài thuyết giáo kéo dài suốt mười lăm phút đồng hồ. Suốt khoảng thời gian đó, nàng Người máy chỉ biết cúi gằm mặt, hai tay chắp lại chịu phạt, miệng liên tục lẩm bẩm lời xin lỗi trong vô vọng.

Phải mất một lúc lâu tôi mới lấy lại được chút bình tĩnh. May cho cô nàng là tôi không có quyền sa thải, nếu không thì Roboco chắc chắn đã phải ra đường đứng rồi. Tôi thở dài, nhìn đống hoang tàn: ``Thôi được rồi, giờ lo mà dọn dẹp đi. Nhưng trước hết, giúp anh nhổ Shion lên khỏi cái hố kia đã.''

Roboco gật đầu lia lịa, tiến đến chỗ Shion đang bị chôn vùi. Cô nàng túm lấy hai tay Shion rồi kéo mạnh, đúng lúc đó nàng Phù thủy bất ngờ bừng tỉnh.

``Ơ\dots\ chuyện gì đang xảy ra thế này!?'' Shion ngơ ngác nhìn xung quanh, thấy mình đang ở trong một cái hố còn đàn chị thì đang cố kéo mình lên. Thế nhưng, dường như có một lực vô hình nào đó ghì chặt lấy em ấy, khiến Roboco dù cố sức đến đâu cũng không thể nhấc nổi, thậm chí còn ngã bật ngửa ra sau.

Roboco la lên thất thanh: ``\textbf{Trời ơi, em ấy như bị dính chặt vào lòng đất luôn rồi! Kuon-senpai cứu em với, em kéo không nổi!}''

Tôi đổ mồ hôi hột: ``Cái gì? Sức mạnh rô-bốt của em mà còn không kéo nổi một đứa bé tí như Shion sao?''

Cuối cùng, tôi đành phải dùng đến giải pháp cuối cùng: gọi điện cho đội cứu hộ và xây dựng chuyên nghiệp đến để đưa Shion lên và đưa em ấy vào viện kiểm tra gấp. Còn Roboco, với danh hiệu ``hậu đậu kiêm phá hoại'' mới được xác lập, bị tôi buộc phải về nhà viết bản kiểm điểm sâu sắc.

\ct

Một tuần sau, Shion được xuất viện sau khi sức khỏe đã hồi phục hoàn toàn. Ngay khi em ấy vừa đặt chân trở lại công ty, tôi đã lập tức giao cho em ấy trọng trách: dùng phép thuật để khôi phục lại tòa nhà văn phòng.

``Anh không cần em phải xây lại cầu kỳ, chỉ cần tái thiết phần cấu trúc cơ bản là được rồi,'' Tôi nói với em ấy, nhưng lòng vẫn đầy băn khoăn. ``Còn phần nội thất\dots\ chắc anh không dám nhờ đâu. Em vừa mới ra viện mà.''

``Được rồi, để em thử xem sao,'' Shion đáp đầy tự tin.

Với quyền năng phù thủy của mình, Shion nhanh chóng dựng lại bộ khung tòa nhà. Tuy nhiên, đúng như tôi dự đoán, phần nội thất là một thảm họa khác. Khi hai chúng tôi bước vào văn phòng mới, đập vào mắt chỉ là bốn bức tường trống huơ trống hoác, chẳng còn lấy một vật dụng nào.

Shion ngơ ngác nhìn quanh: ``Ơ\dots\ cái ghế sofa êm ái của em đâu rồi? Còn bàn làm việc của anh nữa\dots\ Chắc tại lúc nãy em không tập trung vào chi tiết được rồi.''

Tôi thở dài, xoa thái dương: ``Thôi, em mới khỏe lại mà, không trách được. Phần còn lại anh sẽ lo liệu. Tiền mua sắm nội thất mới này sẽ được trích từ lương của anh và tiền phạt của cô nàng Robo-PON kia.''

Shion tròn mắt: ``Tiền phạt của chị ấy? Roboco-senpai chịu trả thật sao?''

Tôi nhún vai đầy vẻ tàn nhẫn: ``Em ấy không có sự lựa chọn nào khác đâu. Không chịu nộp phạt là anh với A-chan sẽ cho nghỉ việc ngay lập tức.''

Shion khẽ cười: ``Em không nghĩ anh chị lại ác thế đâu. Nhưng mà, anh đừng có phạt chị ấy nặng quá nhé. Dù sao thì cũng chỉ là tai nạn ngoài ý muốn thôi mà.''

Tôi nhìn em ấy, giọng trầm xuống: ``Anh biết chứ, anh cũng chẳng muốn đuổi ai cả. Nhưng lần sau, làm ơn trước khi để cô nàng ấy làm điều gì đó `đột phá', hãy suy nghĩ thật kỹ cho anh nhờ.''

Shion bĩu môi: ``Anh cũng có những ý tưởng chả giống ai còn nói người ta! Nào là trò dọa ma Sora-senpai, rồi cả vụ anh đập vợt muỗi vào mặt em nữa, em vẫn chưa quên đâu nhé!''

Tôi hơi ngựng lại, rồi gãi đầu thừa nhận: ``\dots Ừ. Có lẽ anh cũng đã bắt đầu bị lây cái sự `bất bình thường' của mấy đứa rồi cũng nên.''

Đến cuối cùng, nếu phải tính toán thiệt hại thực sự để đền bù cả một tòa nhà, thì cú xoay người của nàng Phù thủy chính là màn nghịch dại tốn kém nhất lịch sử \hll\ từ trước đến nay.

\end{document}
